\chapter{Stochastic Processes, SDEs, and the Ornstein--Uhlenbeck Channel}
\label{ch:stochproc}\label{ch:sde}

Chapter~\ref{ch:statmech} answered the question ``which distribution?'':
Boltzmann's. This chapter answers ``how does probability move?'', in
discrete and in continuous time. Its material enters the thesis in
three places. Markov chains are the \emph{data} whose diffusion score
we compute exactly (the AR(1) sequence of Section~\ref{sec:ar1} is the
running example of every research chapter); stochastic differential
equations are the \emph{corruption and generation processes} of diffusion
models (the Ornstein--Uhlenbeck channel of Section~\ref{sec:ou} appears
on every page of the research chapters); and the Fokker--Planck and
time-reversal theorems of
Sections~\ref{sec:fokker-planck}--\ref{sec:reversal} are the
mathematical basis of the generative programme reviewed in
Chapter~\ref{ch:diffusion-long}. Standard references are
\citet{robert2004monte, brockwell1991time, oksendal2003stochastic,
gardiner2009stochastic, risken1996fokker}.

\section{Stochastic processes and stationarity}\label{sec:stationarity}

A \emph{stochastic process} is a collection of random variables indexed
by time, $(X_k)_{k \geq 0}$ in discrete time or $(X_t)_{t \geq 0}$ in
continuous time, defined on a common probability space; equivalently, a
probability law on the space of trajectories. Two notions of time
homogeneity are used constantly in this thesis and are worth stating
precisely \citep{brockwell1991time}.

\begin{definition}[Strict and weak stationarity]\label{def:stationarity}
A process $(X_k)$ is \emph{strictly stationary} if all its
finite-dimensional laws are invariant under time shifts: for every $n$,
every set of indices $k_1, \dots, k_n$, and every shift $h$,
\begin{equation*}
  (X_{k_1}, \dots, X_{k_n})
  \;\stackrel{d}{=}\;
  (X_{k_1 + h}, \dots, X_{k_n + h}) .
\end{equation*}
It is \emph{weakly (second-order) stationary} if its mean is constant and
its autocovariance depends only on the lag,
$\operatorname{Cov}(X_k, X_{k+d}) = \gamma(d)$.
\end{definition}

Strict stationarity implies weak stationarity whenever second moments
exist; for \emph{Gaussian} processes the two notions coincide, because a
Gaussian law is determined by its first two moments. Stationarity has an
immediate matrix consequence that the research chapters exploit: the
covariance matrix of $K$ consecutive observations of a weakly stationary
process, $(\Sigma)_{ij} = \gamma(i - j)$, is constant along diagonals, a
symmetric \emph{Toeplitz} matrix. Whenever a covariance in this thesis is
Toeplitz, stationarity is the reason.

\section{Markov chains: kernels, stationarity, reversibility}
\label{sec:markov}

\begin{definition}[Markov chain]
A sequence of random variables $(X_k)_{k \geq 0}$ with values in a state
space $\mathcal{X}$ is a \emph{Markov chain} if for every $k$ and every
measurable $A \subseteq \mathcal{X}$,
\begin{equation}
  \Prob\big(X_{k+1} \in A \,\big|\, X_k, X_{k-1}, \dots, X_0\big)
  \;=\; \Prob\big(X_{k+1} \in A \,\big|\, X_k\big) .
  \label{eq:markov-property-long}
\end{equation}
The conditional law is encoded by a \emph{transition kernel}
$M(x' \,|\, x)$: a density (or matrix, for discrete $\mathcal{X}$) with
$\int M(x'|x)\,\dd x' = 1$ for every $x$.
\end{definition}

The Markov property \eqref{eq:markov-property-long} is exactly a statement
about \emph{factorisation}: the joint law of the first $K$ states is
\begin{equation}
  p(x_0, x_1, \dots, x_{K-1})
  \;=\; p_0(x_0)\, \prod_{k=0}^{K-2} M(x_{k+1} \,|\, x_k) .
  \label{eq:chain-factorisation-long}
\end{equation}
This factorisation is the most heavily used equation of the thesis. It says that the joint distribution of a Markov trajectory is a
product of \emph{local} terms, each touching two consecutive variables:
the structure that makes the chain a tree in the factor-graph sense of
Chapter~\ref{ch:graphical}, that makes its precision matrix tridiagonal in
the Gaussian case of Chapter~\ref{ch:gaussian}, and that makes belief
propagation exact.

Marginal distributions evolve by integrating the kernel:
$p_{k+1}(x') = \int M(x'|x)\, p_k(x)\,\dd x$, compactly
$p_{k+1} = M p_k$. A distribution $\pi$ is \emph{stationary} (invariant)
for $M$ if $\pi = M\pi$; a chain started in $\pi$ stays in $\pi$ forever,
and the trajectory is then a strictly stationary process in the sense of
Definition~\ref{def:stationarity}. The kernel satisfies \emph{detailed
balance} with respect to $\pi$ if
\begin{equation}
  \pi(x)\, M(x' \,|\, x) \;=\; \pi(x')\, M(x \,|\, x')
  \qquad \text{for all } x, x' :
  \label{eq:detailed-balance}
\end{equation}
the equilibrium probability flux from $x$ to $x'$ equals the reverse
flux, so the chain, watched in equilibrium, is statistically
indistinguishable run forwards or backwards. Detailed balance implies
stationarity (integrate \eqref{eq:detailed-balance} over $x$), and its
power is practical: it is a local, pointwise condition, checkable without
computing any integral.

Stationarity alone does not guarantee that the chain \emph{finds} $\pi$;
for that one needs the chain to be able to reach everywhere
(\emph{irreducibility}) without getting trapped in cycles
(\emph{aperiodicity}). Under these conditions the fundamental convergence
theorem holds \citep{robert2004monte}: the marginal law of $X_k$
converges to $\pi$ from any initial condition, and time averages converge
to expectations under $\pi$ (the ergodic theorem),
\begin{equation}
  \frac{1}{N} \sum_{k=1}^{N} f(X_k)
  \;\xrightarrow[N \to \infty]{\text{a.s.}}\;
  \E_\pi[f] .
  \label{eq:ergodic}
\end{equation}
Equation \eqref{eq:ergodic} is the rigorous form of the ergodic postulate
invoked in Section~\ref{sec:ensembles}, with the physical dynamics
replaced by an algorithmic one.

\section{The Ornstein--Uhlenbeck process: the corruption channel}
\label{sec:ou}

The noising process used by diffusion models, and by every result of this
thesis, is the \emph{Ornstein--Uhlenbeck process}
\citep{uhlenbeck1930theory}: Brownian motion pulled back to the origin by
a linear spring,
\begin{equation}
  \dd X_t \;=\; -X_t\,\dd t \;+\; \sqrt{2}\;\dd W_t ,
  \qquad X_0 = a .
  \label{eq:ou-sde}
\end{equation}
(The drift and diffusion constants are a normalisation choice, the
variance-preserving convention of the diffusion literature
\citep{song2021scorebased}; any other choice rescales time.) The OU
process is the unique Gaussian, Markov, stationary process in continuous
time, and it is exactly solvable.

\begin{toolbox}{Exact solution of the OU SDE by integrating factor}
Multiply \eqref{eq:ou-sde} by $e^{t}$ and consider $Y_t = e^{t} X_t$. By
It\^o's lemma with $\phi(x,t) = e^{t}x$ (here
$\partial^2 \phi/\partial x^2 = 0$, so no correction term arises):
\begin{equation*}
  \dd Y_t
  = e^{t} X_t\,\dd t + e^{t}\big( -X_t\,\dd t + \sqrt{2}\,\dd W_t \big)
  = \sqrt{2}\, e^{t}\, \dd W_t .
\end{equation*}
The drift cancels exactly; that is what the integrating factor is for.
Integrate from $0$ to $t$ and undo the substitution:
\begin{equation*}
  X_t \;=\; a\, e^{-t} \;+\; \sqrt{2} \int_0^t e^{-(t-s)}\, \dd W_s .
\end{equation*}
The integral is Gaussian (limit of Gaussian sums), with mean zero and
variance given by the It\^o isometry:
\begin{equation*}
  \operatorname{Var}(X_t)
  = 2 \int_0^t e^{-2(t-s)}\, \dd s
  = 2\, e^{-2t}\, \frac{e^{2t} - 1}{2}
  = 1 - e^{-2t} .
\end{equation*}
Hence the \emph{transition kernel} of the OU channel is
\begin{equation*}
  p(x_t \,|\, X_0 = a)
  \;=\; \Nd\!\big( x_t;\; a\,e^{-t},\; \Deltat \big),
  \qquad
  \boxed{\;\Deltat \;=\; 1 - e^{-2t}\;} .
\end{equation*}
\hfill$\square$
\end{toolbox}

The boxed kernel is used on every page of the research chapters, so we
fix its reading now. Corrupting a data point $a$ for a time $t$ produces
\begin{equation}
  x \;=\; a\, e^{-t} \;+\; \sqrt{\Deltat}\;\varepsilon,
  \qquad \varepsilon \sim \Nd(0, 1):
  \label{eq:ou-channel}
\end{equation}
the signal is \emph{contracted} by $e^{-t}$ while isotropic noise of
variance $\Deltat$ fills in. At $t = 0$: $x = a$ (no corruption); as
$t \to \infty$: $x \sim \Nd(0, 1)$ regardless of $a$ (all information
destroyed, the equilibrium of the spring). The signal-to-noise ratio
$e^{-2t}/\Deltat$ sweeps continuously from $\infty$ to $0$, and it
diverges as $1/2t$ for small $t$, a factor that reappears throughout the
thesis as an error amplifier. The OU channel interpolates between the
data distribution and a standard Gaussian, which is precisely what a
diffusion model needs.

Two structural remarks, both load-bearing later. First, the OU kernel is
Gaussian \emph{in $a$ as well as in $x$}: viewed as a function of the
clean variable,
$\Nd(x; a e^{-t}, \Deltat) \propto
\exp[-(e^{-2t}/2\Deltat)(a - x e^{t})^2]$, a Gaussian likelihood factor
of precision $e^{-2t}/\Deltat$ centred at the deconvolved observation
$e^t x$. This is the fact that makes the posteriors of
Chapter~\ref{ch:gaussian} jointly Gaussian and the messages of the BP
computation closable. Second, if a random vector
$a \sim \Nd(\mu, \Sigma_0)$ is pushed through \eqref{eq:ou-channel}
coordinate-wise with \emph{independent} noises, the output is Gaussian
with
\begin{equation}
  \mu_t = e^{-t}\mu,
  \qquad
  \Sigt \;=\; e^{-2t}\, \Sigma_0 \;+\; \Deltat\, I ,
  \label{eq:ou-cov-map}
\end{equation}
since means and covariances transform linearly and the added noise is
white. Equation \eqref{eq:ou-cov-map}, shrink the covariance and add a
multiple of the identity, is the entire forward process of this thesis in
one line.

\section{Time reversal: Anderson's theorem}\label{sec:reversal}

One last piece of SDE theory converts everything above into a generative
method. Suppose data $X_0 \sim p_0$ are corrupted by a forward SDE
\eqref{eq:generic-sde} up to time $T$, producing marginals $p_t$. Can the
corruption be undone? Can we find an SDE which, run from $X_T \sim p_T$
backwards, reproduces the same marginals? The answer
\citep{anderson1982reverse} is yes, and the reverse drift needs exactly
one unknown object:
\begin{equation}
  \dd X_t \;=\;
  \Big[ f(X_t, t) \;-\; g(t)^2\, \nabla_x \log p_t(X_t) \Big]\,\dd t
  \;+\; g(t)\, \dd \bar W_t ,
  \label{eq:reverse-sde}
\end{equation}
where time runs backwards from $T$ to $0$ and $\bar W$ is a reverse-time
Wiener process. The correction to the drift is the \emph{score}
$\nabla_x \log p_t$ of the noised marginal: the object defined in
\eqref{eq:score-final} and computed exactly, for Markov-chain data, in
the research chapters. A heuristic that makes \eqref{eq:reverse-sde}
plausible follows from the Fokker--Planck equation: rewriting its
diffusion term as a transport term,
\begin{equation*}
  \frac{g^2}{2} \frac{\partial^2 p_t}{\partial x^2}
  = \frac{\partial}{\partial x}\!\left(
      \frac{g^2}{2}\, p_t\, \frac{\partial \log p_t}{\partial x}
    \right),
\end{equation*}
one sees that the density evolution of \eqref{eq:generic-sde} is
\emph{identical} to that of the noise-free ODE with velocity
$f - \tfrac{g^2}{2} \nabla \log p_t$ (the \emph{probability-flow ODE} of
\citealp{song2021scorebased}); running that transport backwards, and
re-injecting noise consistently, yields \eqref{eq:reverse-sde}. The
rigorous statement and proof are in \citet{anderson1982reverse}.

\section{From algorithmic chains to data chains}

A closing change of perspective, to set up the research problem. In this
chapter Markov chains began as \emph{algorithms}: their stationary law
was the target, their trajectory a means to an end. From
Chapter~\ref{ch:gaussian} onwards the chain \eqref{eq:ar1-long} plays the
opposite role: it is the \emph{data-generating process}, its trajectory
$(a_0, \dots, a_{K-1})$ is one data point (a video of $K$ frames), and
the object of interest is the probability law of the whole trajectory
after each frame has been independently corrupted by the OU channel
\eqref{eq:ou-channel}. The factorisation \eqref{eq:chain-factorisation-long}
then stops being a computational convenience and becomes the structure to
be discovered: the research programme of this thesis asks how much of
\eqref{eq:chain-factorisation-long} survives noising, and how an inference
algorithm, or a neural network, can exploit what survives. In one
sentence: the OU channel destroys data at a known rate, the
Fokker--Planck equation describes the destruction at the level of
densities, and Anderson's theorem says the destruction is invertible by
anyone who knows the score, which is why the rest of this thesis is about
computing that score exactly when the data have Markov structure.
